% Prajñāpāramitāhṛdaya - Heart Sūtra
% Sanskrit Critical Edition
% Recognizing Chinese Compositional Priority (Nattier 1992)

\documentclass[11pt,a4paper]{article}

% ============================================================================
% PACKAGES
% ============================================================================

\usepackage[margin=2.5cm,bottom=3cm]{geometry}
\usepackage{fontspec}
\usepackage{xeCJK}
\usepackage{xcolor}
\usepackage{longtable}
\usepackage{array}
\usepackage{enumitem}

% Per-page footnotes
\usepackage[perpage]{footmisc}

% Page style
\usepackage{fancyhdr}
\pagestyle{fancy}
\fancyhf{}
\renewcommand{\headrulewidth}{0.4pt}
\fancyhead[L]{\small Prajñāpāramitāhṛdaya -- Sanskrit Critical Edition}
\fancyhead[R]{\small\devanagarifont प्रज्ञापारमिताहृदय}
\fancyfoot[C]{\thepage}

% ============================================================================
% FONTS
% ============================================================================

\setmainfont{Noto Serif}[
  BoldFont = {Noto Serif},
  BoldFeatures = {RawFeature=+axis{wght=700}},
]
\setsansfont{Noto Sans}
\setCJKmainfont{Noto Serif CJK SC}

\newfontfamily\devanagarifont{Noto Serif Devanagari}[
  BoldFont = {Noto Serif Devanagari},
  BoldFeatures = {RawFeature=+axis{wght=700}},
]
\newcommand{\deva}[1]{{\devanagarifont #1}}

% ============================================================================
% CUSTOM COMMANDS
% ============================================================================

\newcommand{\linenum}[1]{\textbf{\small #1.}}
\newcommand{\var}[1]{\textsuperscript{\color{red}\textbf{#1}}}
\newcommand{\lemma}[1]{\textbf{#1}}
\newcommand{\reading}[1]{#1}
\newcommand{\wit}[1]{\textsc{#1}}
\newcommand{\depmark}[2]{{\small\color{blue}[#1\,$\rightarrow$\,#2]}}
\newcommand{\gloss}[1]{{\small\color{teal}\textit{#1}}}
\newcommand{\fluent}[1]{{\small #1}}

% ============================================================================
% DOCUMENT
% ============================================================================

\begin{document}

% ----------------------------------------------------------------------------
% TITLE PAGE
% ----------------------------------------------------------------------------

\begin{center}
{\Huge\bfseries\deva{प्रज्ञापारमिताहृदय}}\\[0.5em]
{\LARGE Prajñāpāramitāhṛdayam}\\[0.3em]
{\Large Heart Sūtra}\\[1.5em]
{\large Sanskrit Critical Edition}\\[0.5em]
{\normalsize Short Recension (after Attwood 2015--2021)}\\[2em]
\end{center}

\hrule
\vspace{1em}

\section*{Editorial Note}

This edition recognizes \textbf{Chinese compositional priority} following Nattier (1992). The Sanskrit text is a back-translation from Chinese, not the original. Accordingly:

\begin{enumerate}
\item The base text follows the earliest Sanskrit witnesses (Hōryūji, 7th--8th c.)
\item Late Tantric additions (oṃ, elaborate maṅgala, na-aprāptiḥ) are relegated to the apparatus
\item Errors from back-translation (e.g., kṣaya for nirodha) are preserved and documented
\item This represents what the Sanskrit translator produced, not a hypothetical ``original''
\end{enumerate}

\section*{Witnesses}

\begin{description}[leftmargin=2cm, labelwidth=1.8cm]
\item[\wit{Ja}] Hōryūji Manuscript (Japan, 7th--8th c.). Earliest Sanskrit witness. Short recension.
\item[\wit{Na}] Nepalese MS Cambridge 1164 (ca. 11th c.). Short recension.
\item[\wit{Nb}] Nepalese MS Cambridge 1165 (17th--18th c.). Long recension with oṃ.
\item[\wit{Nc}] Nepalese MS RAS (19th c.). Late Tantric.
\item[\wit{Nd}] Nepalese MS India Office (18th--19th c.). Extended version.
\item[\wit{Ne}] Nepalese MS NGMPP (15th--17th c.). Long recension with oṃ.
\item[\wit{T256}] Tang Transliteration (8th--9th c.). Sanskrit preserved in Chinese characters.
\end{description}

\newpage

% ----------------------------------------------------------------------------
% TEXT
% ----------------------------------------------------------------------------

\section*{Text}

\begin{center}
{\Large\bfseries Prajñāpāramitāhṛdayam}\\[0.3em]
{\large\deva{प्रज्ञापारमिताहृदयम्}}
\end{center}

\vspace{1em}

\linenum{1} namas sarvajñāya\footnote{\lemma{namas sarvajñāya}] This simple maṅgala appears in \wit{Ja}. Later witnesses have elaborate Tantric formulas: \reading{Oṃ namo Bhagavatyai Ārya-Prajñāpāramitāyai!} \wit{Nb Nc Nd Ne}. The elaboration is a Tantric development (9th--10th c.).} ||

\vspace{0.3em}
\gloss{namas (homage) sarvajñāya (to the omniscient one)}

\vspace{0.2em}
\fluent{Homage to the Omniscient One.}

\vspace{0.5em}

\subsection*{Section I: Opening}

\linenum{2} ārya-avalokiteśvaro\footnote{\lemma{avalokiteśvaro}] \reading{avalokita-} \wit{Ja}. The \wit{Ja} reading may represent earlier confusion. Chinese 觀自在 supports the \textit{-īśvara} compound. \depmark{zh}{sa} The Chinese reading is primary.} bodhisattvo\\
gambhīrāṃ prajñāpāramitā-caryāṃ caramāṇo\\
vyavalokayati sma\footnote{\lemma{vyavalokayati sma}] Unusual periphrastic past tense (\textit{sma} + present). Standard sūtra style uses present or aorist. \depmark{zh}{sa} May reflect Chinese narrative past 時 (\textit{shí}) ``at that time,'' indicating translation from Chinese.}:

\vspace{0.3em}
\linenum{3} pañca-skandhāṃs\footnote{\lemma{pañca-skandhāṃs}] Accusative. Conze (1967) erroneously prints \reading{pañca-skandhāḥ} (nominative). Attwood (2015) corrects: the bodhisattva \emph{sees} the aggregates (accusative required).} tāṃś ca svabhāva-śūnyān paśyati sma ||

\vspace{0.5em}
\gloss{ārya (noble) avalokiteśvaro (Avalokiteśvara) bodhisattvo (bodhisattva), gambhīrāṃ (profound) prajñāpāramitā-caryāṃ (prajñāpāramitā-practice) caramāṇo (practicing), vyavalokayati sma (was looking down): pañca-skandhāṃs (five aggregates) tāṃś ca (and those) svabhāva-śūnyān (empty of intrinsic nature) paśyati sma (he saw).}

\vspace{0.3em}
\fluent{The noble Avalokiteśvara Bodhisattva, while practicing the profound prajñāpāramitā course, looked down and saw the five aggregates, and he saw them as empty of intrinsic nature.}

\vspace{1em}
\subsection*{Section II: Form and Emptiness}

\linenum{4} iha Śāriputra\\
rūpaṃ śūnyatā śūnyataiva rūpam |\\
rūpān na pṛthak śūnyatā\footnote{\lemma{rūpān na pṛthak śūnyatā}] Ablative construction ``emptiness is not separate \emph{from} form.'' Grammatically unusual; may reflect Chinese 色不異空 translated woodenly. \depmark{zh}{sa}} śūnyatāyā na pṛthag rūpam |

\vspace{0.3em}
\linenum{5} yad rūpaṃ sā śūnyatā yā śūnyatā tad rūpam\footnote{\lemma{yad rūpaṃ sā śūnyatā...}] This line (``that which is form, that is emptiness...'') has \emph{no Chinese parallel}. It appears to be an addition made during Sanskrit composition to clarify the form-emptiness equation. Absent from T251, Fangshan stele, and all Chinese witnesses.} |

\vspace{0.3em}
\linenum{6} evam eva vedanā-saṃjñā-saṃskāra-vijñānāni ||

\vspace{0.5em}
\gloss{iha (here) Śāriputra: rūpaṃ (form) śūnyatā (emptiness), śūnyataiva (emptiness itself) rūpam (form); rūpān (from form) na pṛthak (not separate) śūnyatā (emptiness), śūnyatāyā (from emptiness) na pṛthag (not separate) rūpam (form); yad (which) rūpaṃ (form) sā (that) śūnyatā (emptiness), yā (which) śūnyatā (emptiness) tad (that) rūpam (form). evam eva (just so) vedanā (feeling) saṃjñā (perception) saṃskāra (formation) vijñānāni (consciousnesses).}

\vspace{0.3em}
\fluent{Here, Śāriputra, form is emptiness, and emptiness itself is form. Emptiness is not separate from form; form is not separate from emptiness. That which is form is emptiness; that which is emptiness is form. Just so are feeling, perception, formation, and consciousness.}

\vspace{1em}
\subsection*{Section III: Characteristics of Emptiness}

\linenum{7} iha Śāriputra sarva-dharmāḥ śūnyatā-lakṣaṇā |\\
anutpannā aniruddhā |\\
amalā avimalā |\\
anūnā aparipūrṇāḥ ||

\vspace{0.5em}
\gloss{iha (here) Śāriputra: sarva-dharmāḥ (all dharmas) śūnyatā-lakṣaṇā (emptiness-marked); anutpannā (not arisen) aniruddhā (not ceased); amalā (not defiled) avimalā (not undefiled); anūnā (not deficient) aparipūrṇāḥ (not complete).}

\vspace{0.3em}
\fluent{Here, Śāriputra, all dharmas are marked by emptiness: not arisen, not ceased; not defiled, not undefiled; not deficient, not complete.}

\vspace{1em}
\subsection*{Section IV: Negation Series}

\linenum{8} tasmāc Chāriputra śūnyatāyāṃ\\
na rūpaṃ na vedanā na saṃjñā na saṃskārāḥ na vijñānam |

\vspace{0.3em}
\linenum{9} na cakṣuḥ-śrotra-ghrāṇa-jihvā-kāya-manāṃsi |\\
na rūpa-śabda-gandha-rasa-spraṣṭavya-dharmāḥ |

\vspace{0.3em}
\linenum{10} na cakṣur-dhātur yāvan na manovijñāna-dhātuḥ |

\vspace{0.3em}
\linenum{11} na-avidyā na-avidyā-kṣayo\footnote{\lemma{avidyā-kṣayo... jarā-maraṇa-kṣayaḥ}] \textbf{Key evidence for back-translation.} The Sanskrit uses \textit{kṣaya} (``destruction, wearing away'') where the Large Prajñāpāramitā (Sanskrit Pañcaviṃśatisāhasrikā) uses \textit{nirodha} (``cessation''). Chinese 盡 (\textit{jìn}) can translate either term. A translator working \emph{from} Sanskrit would have used \textit{nirodha} (standard Prajñāpāramitā terminology). The use of \textit{kṣaya} indicates someone working \emph{from} Chinese chose an approximate synonym. This is Nattier's ``smoking gun'' (1992:175--178). \depmark{zh}{sa}}\\
yāvan na jarā-maraṇaṃ na jarā-maraṇa-kṣayaḥ |

\vspace{0.3em}
\linenum{12} na duḥkha-samudaya-nirodha-mārgāḥ |\\
na jñānaṃ na prāptiḥ\footnote{\lemma{na prāptiḥ}] Late witnesses add \reading{na-aprāptiḥ} (``no non-attainment'') \wit{T256 Nb Nc Nd Ne}. This is absent from \wit{T251}, \wit{Fangshan} (661 CE), and \wit{Ja} (7th--8th c.). It is an interpolation in the Sanskrit tradition, subsequently translated into T257 as 亦無無得.} ||

\vspace{0.5em}
\gloss{tasmāc (therefore) Chāriputra (Śāriputra): śūnyatāyāṃ (in emptiness) na rūpaṃ (no form) na vedanā (no feeling) na saṃjñā (no perception) na saṃskārāḥ (no formations) na vijñānam (no consciousness); na cakṣuḥ (no eye) śrotra (ear) ghrāṇa (nose) jihvā (tongue) kāya (body) manāṃsi (minds); na rūpa (no form) śabda (sound) gandha (smell) rasa (taste) spraṣṭavya (tangible) dharmāḥ (mental objects); na cakṣur-dhātur (no eye-element) yāvan (up to) na manovijñāna-dhātuḥ (no mind-consciousness element). na-avidyā (no ignorance) na-avidyā-kṣayo (no destruction of ignorance), yāvan (up to) na jarā-maraṇaṃ (no old age and death) na jarā-maraṇa-kṣayaḥ (no destruction of old age and death). na duḥkha (no suffering) samudaya (origination) nirodha (cessation) mārgāḥ (paths); na jñānaṃ (no knowledge) na prāptiḥ (no attainment).}

\vspace{0.3em}
\fluent{Therefore, Śāriputra, in emptiness there is no form, no feeling, no perception, no formations, no consciousness; no eye, ear, nose, tongue, body, or mind; no form, sound, smell, taste, tangible, or mental object; no eye-element, up to and including no mind-consciousness element. No ignorance, no destruction of ignorance, up to and including no old age and death, no destruction of old age and death. No suffering, origination, cessation, or path; no knowledge, no attainment.}

\vspace{1em}
\subsection*{Section V: Bodhisattva's Result}

\linenum{13} tasmāc Chāriputra aprāptitvād\\
bodhisattvasya prajñāpāramitām āśritya\\
viharaty acittāvaraṇaḥ |

\vspace{0.3em}
\linenum{14} cittāvaraṇa-nāstitvād atrastro\\
viparyāsa-atikrānto\\
niṣṭhā-nirvāṇaḥ\footnote{\lemma{niṣṭhā-nirvāṇaḥ}] Conze prints \reading{niṣṭhā-nirvāṇa-prāptaḥ}. Attwood (2018) argues this compound is problematic. The Chinese 究竟涅槃 (\textit{jiūjìng nièpán}) ``ultimate nirvāṇa'' is clear; the Sanskrit may be a calque. The \textit{-prāptaḥ} (``attained'') appears to be a later addition for grammatical clarity.} ||

\vspace{0.5em}
\gloss{tasmāc (therefore) Chāriputra (Śāriputra): aprāptitvād (because of non-attainment) bodhisattvasya (of the bodhisattva) prajñāpāramitām (prajñāpāramitā) āśritya (relying on) viharaty (dwells) acittāvaraṇaḥ (with mind unobstructed); cittāvaraṇa-nāstitvād (because of the absence of mental obstruction) atrastro (unafraid) viparyāsa-atikrānto (having gone beyond error) niṣṭhā-nirvāṇaḥ (one whose nirvāṇa is complete).}

\vspace{0.3em}
\fluent{Therefore, Śāriputra, because of non-attainment, the bodhisattva, relying on prajñāpāramitā, dwells with mind unobstructed. Because of the absence of mental obstruction, he is unafraid, has gone beyond error, and his nirvāṇa is complete.}

\vspace{1em}
\subsection*{Section VI: Buddha's Result}

\linenum{15} tryadhva-vyavasthitāḥ sarva-buddhāḥ\\
prajñāpāramitām āśritya\\
anuttarāṃ samyak-saṃbodhim abhisaṃbuddhāḥ ||

\vspace{0.5em}
\gloss{tryadhva-vyavasthitāḥ (abiding in the three times) sarva-buddhāḥ (all buddhas) prajñāpāramitām (prajñāpāramitā) āśritya (relying on) anuttarāṃ (unsurpassed) samyak-saṃbodhim (perfect complete enlightenment) abhisaṃbuddhāḥ (fully awakened to).}

\vspace{0.3em}
\fluent{All buddhas abiding in the three times, relying on prajñāpāramitā, have fully awakened to unsurpassed, perfect, complete enlightenment.}

\vspace{1em}
\subsection*{Section VII: Mantra Praise}

\linenum{16} tasmāj jñātavyam:\\
prajñāpāramitā mahā-mantro\\
mahā-vidyā-mantro\\
'nuttara-mantro\\
'samasama-mantraḥ |

\vspace{0.3em}
\linenum{17} sarva-duḥkha-praśamanaḥ |\\
satyam amithyatvāt ||

\vspace{0.5em}
\gloss{tasmāj (therefore) jñātavyam (it should be known): prajñāpāramitā (perfection of wisdom) mahā-mantro (great mantra) mahā-vidyā-mantro (mantra of great knowledge) 'nuttara-mantro (unsurpassed mantra) 'samasama-mantraḥ (mantra equal to the unequalled); sarva-duḥkha-praśamanaḥ (allayer of all suffering); satyam (true) amithyatvāt (because not false).}

\vspace{0.3em}
\fluent{Therefore it should be known: prajñāpāramitā is the great mantra, the mantra of great knowledge, the unsurpassed mantra, the mantra equal to the unequalled, which allays all suffering --- true, because not false.}

\vspace{1em}
\subsection*{Section VIII: Mantra}

\linenum{18} prajñāpāramitāyām ukto mantraḥ | tadyathā:

\vspace{0.5em}
\begin{center}
{\Large gate gate pāragate pārasaṃgate bodhi svāhā\footnote{\lemma{gate gate...}] Late witnesses prefix \reading{oṃ} \wit{Nb Ne T257}. This is a Tantric addition (9th--10th c.), absent from \wit{T251}, \wit{Fangshan} (661 CE), \wit{Ja}, \wit{Na}, and \wit{T256}. All early witnesses begin directly with \textit{gate}.} ||}

\vspace{0.3em}
{\large\deva{गते गते पारगते पारसंगते बोधि स्वाहा}}
\end{center}

\vspace{0.5em}
\linenum{19} iti prajñāpāramitā-hṛdayaṃ samāptam ||

\vspace{0.5em}
\gloss{prajñāpāramitāyām (in the prajñāpāramitā) ukto (proclaimed) mantraḥ (mantra); tadyathā (it runs thus): gate gate (gone, gone) pāragate (gone beyond) pārasaṃgate (gone completely beyond) bodhi (awakening) svāhā (hail!). iti (thus) prajñāpāramitā-hṛdayaṃ (heart of the perfection of wisdom) samāptam (is complete).}

\vspace{0.3em}
\fluent{The mantra is proclaimed in the prajñāpāramitā. It runs thus: Gone, gone, gone beyond, gone completely beyond --- awakening! Hail! Thus concludes the Heart of the Perfection of Wisdom.}

\newpage

% ----------------------------------------------------------------------------
% NOTE ON CONZE
% ----------------------------------------------------------------------------

\section*{Note on Conze's Edition}

Edward Conze's critical editions (1948, 1967, 1975) applied classical stemmatic method, treating Sanskrit as the original language and privileging late Nepalese manuscripts. This produced a conflated text containing:

\begin{itemize}
\item Late Tantric additions (oṃ, elaborate maṅgala) presented as original
\item Interpolations (na-aprāptiḥ) included in the main text
\item Grammatical ``corrections'' (nominative for accusative in Section I)
\item The compound \textit{niṣṭhā-nirvāṇa-prāptaḥ} (questioned by Attwood)
\end{itemize}

Attwood (2015--2024) has systematically identified these problems and proposed corrections. His revised edition (Attwood 2024) presents an emended Sanskrit text with nine specific corrections, accompanied by Devanāgarī, IAST romanization, and an English translation based on his ``phenomenological reading'' of the text. Attwood's corrections include the accusative \textit{pañcaskandhāṃs} (for Conze's nominative), removal of the \textit{na-aprāptiḥ} interpolation, omission of the Tantric \textit{oṃ} and \textit{maṅgala}, and the \textit{yad rūpaṃ} line (no Chinese counterpart).

This edition follows Attwood in making the same substantive corrections but differs in format: where Attwood presents a single revised text, we present a \emph{critical} edition with formal witness sigla and an apparatus recording manuscript variants. Our footnotes explain the Chinese-priority reasoning behind each editorial decision, allowing readers to evaluate the evidence independently.

\vspace{2em}
\hrule
\vspace{0.5em}
\begin{center}
\textit{Sanskrit Critical Edition} $\cdot$ Prajñāpāramitāhṛdayam $\cdot$ Recognizing Chinese Priority
\end{center}

\end{document}
